\paragraph{3. The integral $j_*$ condition survives ramified transfer.}
For a number field $K$, let $X_K=\operatorname{Spec}\mathcal O_K[1/2]$
and $j_K:\operatorname{Spec}K\to X_K$. At finite coefficients
$M=T/2^aT$, the low-degree Leray sequence gives an injection
\[
 H^1(X_K,j_{K,*}M)\longrightarrow H^1(K,M)
\]
whose image is the kernel of all inertia restrictions outside two.
Indeed the higher-direct-image stalk formula and Leray sequence
\cite[Tags 03Q9 and 03QA]{stackscohom} identify the stalks of
$R^1j_{K,*}M$ at a finite place with $H^1(I_v,M)$;
the edge map is restriction to these stalks. The generic-point map
$j_K$ is affine, hence quasi-compact and quasi-separated, as required
by the stalk formula. Kato's section~8.2,
pp.~180--181, defines the integral groups by the inverse limit of
these finite-coefficient groups, so this assertion applies to the
classes (K4) before rationalization.

Let $L/K$ be finite, with absolute Galois groups $H\subset G_K$.
For a class $z\in H^1(H,M)$ and an inertia subgroup $I\subset G_K$,
the Mackey formula is
\[
 \operatorname{res}_{I}^{G_K}\operatorname{cor}_{H}^{G_K}z
 =\sum_{g\in I\backslash G_K/H}
   \operatorname{cor}_{I\cap gHg^{-1}}^{I}
   \operatorname{res}_{I\cap gHg^{-1}}^{gHg^{-1}}(g z).    \tag{K7}
\]
The intersection in each summand is the upper inertia subgroup
for the corresponding place of $L$. If all upper inertia
restrictions vanish, every summand is zero. This does not require
$L/K$ to be unramified. It is a finite double-coset identity at each
finite coefficient precision. Uniqueness under the Leray injection
makes the extensions across $X_K$ compatible in $a$. Thus the
integral $j_*$ condition survives corestriction, including at primes
ramified in a quadratic twisting field.

